% Overleaf-ready draft for revising Section II/III organization.
% Recommended use:
% 1) Replace the current Fig. 1 caption.
% 2) Replace the current Section II with the block below.
% 3) Replace the current Section III organization with the block below.
% 4) In Section III-B, replace "channel-wise KL divergence" with
%    "dimension-wise KL divergence" for terminology consistency.

% ------------------------------------------------------------------
% Fig. 1 caption
% ------------------------------------------------------------------
% \caption{Generic split federated learning framework considered in this work. Multiple clients collaboratively train a split model by transmitting smashed data to the server over wireless uplink channels and receiving intermediate gradients for backpropagation.}
% \label{fig:system_model}
%
% Optional Fig. 3 caption:
% \caption{Proposed CA-SSFL semantic encoder with VIB-driven semantic masking, SNR-driven channel masking, and dynamic slicing.}
% \label{fig:proposed_encoder}

% ------------------------------------------------------------------
% Section II replacement
% ------------------------------------------------------------------
\section{System Model}

We consider an SFL framework that combines split learning and federated aggregation \cite{gupta2018distributed,thapa2022splitfed}, as illustrated in Fig.~\ref{fig:system_model}. Let $\mathcal{I}=\{1,\ldots,I\}$ denote the client set, and let $\mathcal{D}_i$ denote the local dataset at client $i \in \mathcal{I}$. In the considered setting, clients communicate with the server over wireless uplink channels. The global model is partitioned into a client-side sub-model $\phi_c$ and a server-side sub-model $\phi_s$. For an input sample $x$, the client first computes an intermediate feature representation, commonly referred to as smashed data, as
\begin{equation}
f = \phi_c(x).
\end{equation}
\noindent
The smashed data $f$ is then transmitted to the server through the wireless uplink. The transmitted representation is corrupted by channel noise and fading during uplink transmission, and the received representation at the server can be modeled as
\begin{equation}
f' = h \cdot f + n,
\end{equation}
where $n \sim \mathcal{CN}(0,\sigma_n^2 I)$ denotes additive white Gaussian noise and $h$ denotes the channel coefficient. In particular, $h=1$ for the AWGN channel and $h \sim \mathcal{CN}(0,1)$ for the Rayleigh fading channel. Based on the received representation, the server performs the remaining layers of the split model to produce the final prediction
\begin{equation}
\hat{y} = \phi_s(f').
\end{equation}
This split architecture reduces the computational burden on client devices by offloading most of the deep neural network computation to the server, which makes it attractive for resource-limited IoT environments. In the generic SFL setting, the split model is trained end-to-end using a standard supervised objective defined on the server output. For a ground-truth label $y$, the task loss is written as
\begin{equation}
\mathcal{L}_{\mathrm{task}} = \ell(\hat{y}, y),
\end{equation}
where $\ell(\cdot,\cdot)$ denotes a classification loss such as cross-entropy. During training, the server computes the loss from the final prediction and propagates the intermediate gradient back through the split point so that both $\phi_s$ and $\phi_c$ are jointly optimized. After local updates, the client-side parameters are periodically aggregated across participating clients in a federated manner to form the global model for the next communication round.

Although SFL is effective in reducing on-device computation, its communication cost remains substantial because high-dimensional smashed data must be transmitted at every training iteration \cite{gao2020end}. This problem becomes more severe in wireless environments, where limited bandwidth and channel noise can significantly degrade transmission efficiency and downstream task performance. These challenges motivate a communication-efficient and channel-adaptive transmission mechanism, which is introduced in the next section.

% ------------------------------------------------------------------
% Section III organization replacement
% ------------------------------------------------------------------
\section{Proposed Method}

\subsection{Overview of CA-SSFL}
Building on the SFL framework in Section II, we propose a channel-aware semantic split federated learning (CA-SSFL) architecture to reduce communication overhead while preserving task-relevant information under dynamic channel conditions. Motivated by task-oriented semantic communication \cite{qin2021semantic,gunduz2022beyond}, the proposed method augments the conventional client-server split model with a semantic encoder at the client side and a semantic decoder at the server side, as illustrated in Fig.~\ref{fig:proposed_encoder}. Accordingly, the client-side model is defined as $\Theta_C = \{\phi_c, \theta_{en}\}$, where $\phi_c$ extracts smashed data and $\theta_{en}$ converts it into a transmission-efficient latent representation. The server-side model is defined as $\Theta_S = \{\theta_{de}, \phi_s\}$, where $\theta_{de}$ reconstructs the received latent representation and $\phi_s$ performs the downstream prediction task. Therefore, the overall model is given by $\Theta = \{\phi_c, \theta_{en}, \theta_{de}, \phi_s\}$.

Unlike conventional SFL, the proposed CA-SSFL incorporates two key mechanisms in the client-side semantic encoder: 1) a VIB-based semantic masking module that estimates the task relevance of latent features, and 2) an SNR-driven channel masking module that determines their transmission feasibility under the current wireless condition. By combining these two masks, the encoder performs dynamic slicing and transmits only the feature dimensions that are both semantically important and channel-feasible. The server then restores the sparse latent representation to its original dimensionality and processes it through the semantic decoder and the remaining server-side network.

\subsection{Semantic Encoder and VIB-Driven Semantic Masking}
Let $f_{\mathrm{flt}} \in \mathbb{R}^{K}$ denote the flattened smashed data extracted from the client-side model. The semantic encoder first maps $f_{\mathrm{flt}}$ to the latent statistics of a variational representation, namely the mean $\mu$ and variance $\sigma^2$. To enable stochastic training, we adopt the reparameterization trick \cite{kingma2013auto} and sample the latent vector as
\begin{equation}
z = \mu + \sigma \odot \epsilon, \quad \epsilon \sim \mathcal{N}(0, I).
\end{equation}
Rather than treating all latent dimensions equally, CA-SSFL evaluates their semantic importance through dimension-wise KL divergence between the posterior and a standard normal prior following the variational information bottleneck principle \cite{alemi2017deep}. In implementation, the latent statistics used for KL evaluation are modulated by the current channel condition through the channel-aware encoder path. Therefore, the semantic importance is computed from channel-modulated latent statistics rather than from a strictly channel-independent branch. Based on this score, dimensions with KL values larger than a predefined semantic threshold $\tau_{\mathrm{VIB}}$ are retained, yielding the binary semantic mask $M_{\mathrm{vib}}$.

\subsection{SNR-Driven Channel Masking Module}
To reflect the instantaneous wireless condition, the normalized SNR $s_{\mathrm{norm}} \in [0,1]$ is fed into a lightweight masking network that generates a channel-feasibility score for each latent dimension. Specifically, the masking path produces a probability vector $P \in [0,1]^K$, where each element indicates how likely the corresponding latent dimension is to be reliably transmitted under the current channel state. To adapt the transmission budget, we introduce a dynamic threshold
\begin{equation}
\tau_{\mathrm{ch}} = \tau_{\max} - s_{\mathrm{norm}}(\tau_{\max}-\tau_{\min}).
\end{equation}
Under low-SNR conditions, $\tau_{\mathrm{ch}}$ approaches $\tau_{\max}$, which enforces stronger compression. Under high-SNR conditions, it approaches $\tau_{\min}$, allowing more semantic features to be transmitted. The binary channel mask $M_{\mathrm{ch}}$ is then obtained by comparing $P$ with $\tau_{\mathrm{ch}}$, so that this module serves as a channel-dependent feasibility gate.

\subsection{Dual-Masking and Dynamic Slicing}
The final transmission support is determined by the element-wise product of the semantic mask and the channel mask,
\begin{equation}
M_{\mathrm{final}} = M_{\mathrm{vib}} \odot M_{\mathrm{ch}}.
\end{equation}
The sampled latent vector is masked as $z_{\mathrm{masked}} = z \odot M_{\mathrm{final}}$. To realize actual communication reduction, the client transmits only the active latent values and their indices instead of sending the dense masked vector. At the server, the sparse latent representation is restored to the original dimensionality through zero-filling before being processed by the semantic decoder. Since extreme sparsity can weaken the backpropagated learning signal, we additionally apply sparsity-aware gradient compensation. Let $p$ denote the active feature ratio. During backpropagation, the restored gradient is scaled by $1/\max(p, 0.01)$ so that the effective gradient magnitude remains stable even when the transmission budget varies significantly, following the intuition of inverted dropout scaling \cite{gradientScale}.

\subsection{Training Objective}
The end-to-end training objective of CA-SSFL combines the downstream task loss with a VIB-based compression penalty. Specifically, we minimize
\begin{equation}
\mathcal{L}_{\mathrm{total}} = \mathcal{L}_{\mathrm{task}} + \beta \mathcal{L}_{\mathrm{comp}},
\end{equation}
where $\beta$ controls the trade-off between task performance and information compression. The first term encourages accurate downstream prediction after semantic transmission through the noisy wireless channel, while the second term penalizes redundant latent information through KL regularization from an information bottleneck perspective \cite{alemi2017deep,wang2023performance}. Unlike the generic SFL objective in Section II, this proposed objective explicitly learns transmission-efficient latent features jointly with semantic masking, channel masking, and semantic decoding. As a result, the encoder adaptively allocates the transmitted support according to both semantic relevance and instantaneous channel conditions.

% Suggested Section III table of contents:
% A. Overview of CA-SSFL
% B. Semantic Encoder and VIB-Driven Semantic Masking
% C. SNR-Driven Channel Masking Module
% D. Dual-Masking and Dynamic Slicing
% E. Training Objective

% ------------------------------------------------------------------
% Terminology fix for Section III-B
% ------------------------------------------------------------------
% Replace:
% "we utilize the channel-wise KL divergence as a metric for information significance"
%
% With:
% "we utilize the dimension-wise KL divergence as a metric for information significance"
